\documentclass[11pt,a4paper]{article}

\usepackage[a4paper, margin=2.5cm]{geometry}
\usepackage{booktabs}
\usepackage{array}
\usepackage{longtable}
\usepackage{hyperref}
\usepackage{xcolor}
\usepackage{titlesec}
\usepackage{parskip}
\usepackage{fancyhdr}
\usepackage{listings}

\definecolor{codegray}{rgb}{0.95,0.95,0.95}
\lstset{
  backgroundcolor=\color{codegray},
  basicstyle=\ttfamily\small,
  frame=single,
  framesep=4pt,
  breaklines=true
}

\hypersetup{colorlinks=true, linkcolor=blue, urlcolor=blue}

\pagestyle{fancy}
\fancyhf{}
\rhead{CIF Microarchitecture}
\lhead{Read Transaction Table}
\rfoot{\thepage}

\title{\textbf{CIF Microarchitecture Specification}\\[0.4em]
\large Block: Read Transaction Table (RTT)}
\author{CIF Design Team}
\date{\today}

\begin{document}

\maketitle
\thispagestyle{fancy}
\tableofcontents
\newpage

% -----------------------------------------------------------------------
\section{Overview}

The Read Transaction Table (RTT) is a per-transaction tracking structure
for the read path. It is allocated at tag assignment time and tracks
completion progress for every outstanding read transaction. When all
expected packets for a transaction have been received from MC Core,
the RTT asserts \texttt{rob\_last}, signalling the final beat to the
ROB and Merge Logic. The entry is freed in the same cycle.

The RTT operates independently of the ROB — both blocks tap the same
MC Core completion bus, but serve different purposes: the ROB stores
and reorders data; the RTT tracks transaction-level progress and drives
\texttt{rob\_last}.

% -----------------------------------------------------------------------
\section{Design Parameters}

\begin{table}[h!]
\centering
\begin{tabular}{lll}
\toprule
\textbf{Parameter} & \textbf{Value} & \textbf{Description} \\
\midrule
\texttt{AXID\_W}         & 6   & AXI ID width \\
\texttt{SEQNUM\_W}       & 2   & Seqnum width \\
\texttt{TXN\_TAG\_W}     & 8   & Tag width = AXID\_W + SEQNUM\_W \\
\texttt{RTT\_DEPTH}      & 256 & Entries = $2^{\texttt{AXID\_W} + \texttt{SEQNUM\_W}}$ \\
\texttt{EXP\_PKTS\_W}    & 8   & Expected packets width (covers ARLEN+1, max 256) \\
\texttt{CMP\_PKTS\_W}    & 8   & Completed packets counter width \\
\bottomrule
\end{tabular}
\caption{RTT parameters}
\end{table}

% -----------------------------------------------------------------------
\section{Internal Blocks}

\subsection{RTT Register File}

\begin{itemize}
  \item \textbf{Structure:} 256-entry register file. Each entry stores:
  \begin{itemize}
    \item \texttt{AXID[5:0]} — AXI ID of the transaction
    \item \texttt{exp\_pkts[7:0]} — expected packet count (= \texttt{ARLEN+1})
    \item \texttt{cmp\_pkts[7:0]} — completed packet count (incremented on each MC Core completion)
  \end{itemize}
  \item \textbf{Addressing:} \texttt{addr[7:0] = TXN\_TAG[7:0] = \{AXID[5:0], seqnum[1:0]\}}.
        The tag directly indexes the register file — no translation required.
  \item \textbf{Write ports:}
  \begin{itemize}
    \item Alloc Write Logic writes \texttt{exp\_pkts} on allocation.
    \item Completion Update Logic increments \texttt{cmp\_pkts} on each
          \texttt{RD\_DATA} completion.
    \item Free Logic clears the entry (zeroes \texttt{cmp\_pkts}, invalidates)
          on \texttt{rob\_last}.
  \end{itemize}
  \item \textbf{Read port:} Continuously exposes \texttt{cmp\_pkts} and
        \texttt{exp\_pkts} to Compare Logic.
\end{itemize}

\subsection{Alloc Write Logic}

\begin{itemize}
  \item Activated on \texttt{TAG\_VALID} from the Tag Allocator.
  \item Receives \texttt{TXN\_TAG[7:0]} and \texttt{ARLEN[7:0]}.
  \item Computes \texttt{exp\_pkts = ARLEN + 1}.
  \item Issues \texttt{write\_en} to the RTT Register File at address
        \texttt{TXN\_TAG[7:0]}, writing \texttt{exp\_pkts} and
        initialising \texttt{cmp\_pkts = 0}.
\end{itemize}

\subsection{Completion Update Logic}

\begin{itemize}
  \item Taps the MC Core completion bus: \texttt{RESP\_VALID},
        \texttt{resp\_type[1:0]}, \texttt{tag[7:0]}.
  \item \textbf{Filter:} acts only when \texttt{resp\_type == 2'b00}
        (RD\_DATA). WR\_ACK completions (\texttt{resp\_type == 2'b01})
        are destined for the Response Tracker and are ignored. ERR
        completions for read transactions (\texttt{resp\_type == 2'b10})
        are treated as a valid completion event — \texttt{cmp\_pkts} is
        still incremented so the transaction reaches \texttt{rob\_last}
        and completes normally; error status propagation is handled by
        the ROB and R Response Generator.
  \item Asserts \texttt{inc\_en} to the RTT Register File at address
        \texttt{tag[7:0]}, incrementing \texttt{cmp\_pkts} by 1.
\end{itemize}

\subsection{Compare Logic}

\begin{itemize}
  \item Purely combinational.
  \item Reads \texttt{cmp\_pkts[7:0]} and \texttt{exp\_pkts[7:0]} from
        the RTT Register File for the entry currently being updated.
  \item Asserts \texttt{rob\_last} as a one-cycle pulse when
        \texttt{cmp\_pkts == exp\_pkts}.
  \item \texttt{rob\_last} is forwarded to two destinations
        simultaneously: Free Logic (to clear the entry) and the ROB
        (passed through to Merge Logic coincident with the final data
        beat).
\end{itemize}

\subsection{Free Logic}

\begin{itemize}
  \item Activated by \texttt{rob\_last} from Compare Logic.
  \item Asserts \texttt{clear\_en} to the RTT Register File at address
        \texttt{tag[7:0]}, clearing the entry in the same cycle as
        \texttt{rob\_last}.
  \item Simultaneously asserts \texttt{RD\_FREE\_VALID} and drives
        \texttt{RD\_FREE\_TAG[7:0]} to the Tag Allocator's RD\_FREE
        port, triggering occupancy decrement for the freed AXID.
\end{itemize}

% -----------------------------------------------------------------------
\section{Interface Signals}

\subsection{Inputs}

\begin{longtable}{p{3cm}p{1cm}p{3cm}p{7cm}}
\toprule
\textbf{Signal} & \textbf{Width} & \textbf{Source} & \textbf{Description} \\
\midrule
\texttt{TXN\_TAG[7:0]}    & 8 & Tag Allocator & Allocated tag; indexes the register file \\
\texttt{ARLEN[7:0]}       & 8 & Tag Allocator & Burst length; \texttt{exp\_pkts = ARLEN+1} \\
\texttt{TAG\_VALID}       & 1 & Tag Allocator & Allocation strobe \\
\texttt{RESP\_VALID}      & 1 & MC Core       & Completion valid \\
\texttt{resp\_type[1:0]}  & 2 & MC Core       & 00=RD\_DATA, 01=WR\_ACK, 10=ERR; RTT acts on RD\_DATA and ERR \\
\texttt{tag[7:0]}         & 8 & MC Core       & Completion tag; used as register file address \\
\bottomrule
\caption{RTT input signals}
\end{longtable}

\subsection{Outputs}

\begin{longtable}{p{3cm}p{1cm}p{3cm}p{7cm}}
\toprule
\textbf{Signal} & \textbf{Width} & \textbf{Destination} & \textbf{Description} \\
\midrule
\texttt{rob\_last}          & 1 & ROB / Merge Logic             & Pulse on final packet of a transaction \\
\texttt{RD\_FREE\_VALID}    & 1 & Tag Allocator (RD\_FREE port) & Free strobe; coincident with \texttt{rob\_last} \\
\texttt{RD\_FREE\_TAG[7:0]} & 8 & Tag Allocator (RD\_FREE port) & Tag being freed \\
\bottomrule
\caption{RTT output signals}
\end{longtable}

% -----------------------------------------------------------------------
\section{Internal Signal Table}

\begin{longtable}{p{2.8cm}p{1cm}p{2.5cm}p{2.5cm}p{5.5cm}}
\toprule
\textbf{Signal} & \textbf{Width} & \textbf{From} & \textbf{To} & \textbf{Description} \\
\midrule
\texttt{write\_en}        & 1   & Alloc Write Logic       & RTT Reg File         & Write strobe on allocation \\
\texttt{wr\_addr[7:0]}    & 8   & Alloc Write Logic       & RTT Reg File         & = \texttt{TXN\_TAG[7:0]} \\
\texttt{exp\_pkts[7:0]}   & 8   & Alloc Write Logic       & RTT Reg File         & = \texttt{ARLEN+1} \\
\texttt{inc\_en}          & 1   & Completion Update Logic & RTT Reg File         & Increment \texttt{cmp\_pkts} strobe \\
\texttt{inc\_addr[7:0]}   & 8   & Completion Update Logic & RTT Reg File         & = \texttt{tag[7:0]} from MC Core \\
\texttt{cmp\_pkts[7:0]}   & 8   & RTT Reg File            & Compare Logic        & Current completed count \\
\texttt{exp\_pkts[7:0]}   & 8   & RTT Reg File            & Compare Logic        & Expected count (read port) \\
\texttt{rob\_last}        & 1   & Compare Logic           & Free Logic, ROB      & One-cycle pulse on final completion \\
\texttt{clear\_en}        & 1   & Free Logic              & RTT Reg File         & Clear entry strobe \\
\texttt{clear\_addr[7:0]} & 8   & Free Logic              & RTT Reg File         & = \texttt{tag[7:0]} of completing txn \\
\bottomrule
\caption{RTT internal signals}
\end{longtable}

% -----------------------------------------------------------------------
\section{Functional Behaviour}

\subsection{Allocation Flow}

\begin{enumerate}
  \item Tag Allocator asserts \texttt{TAG\_VALID} with \texttt{TXN\_TAG[7:0]}
        and \texttt{ARLEN[7:0]}.
  \item Alloc Write Logic computes \texttt{exp\_pkts = ARLEN + 1}.
  \item RTT Register File is written at address \texttt{TXN\_TAG[7:0]}:
        \texttt{exp\_pkts} is stored, \texttt{cmp\_pkts} is initialised
        to zero.
\end{enumerate}

\subsection{Completion Flow}

\begin{enumerate}
  \item MC Core asserts \texttt{RESP\_VALID} with \texttt{resp\_type == 2'b00}
        (RD\_DATA) and \texttt{tag[7:0]}.
  \item Completion Update Logic asserts \texttt{inc\_en} at address
        \texttt{tag[7:0]}: \texttt{cmp\_pkts} is incremented by 1.
  \item Compare Logic reads the updated \texttt{cmp\_pkts} and
        \texttt{exp\_pkts} combinationally.
  \item If \texttt{cmp\_pkts == exp\_pkts}: \texttt{rob\_last} is
        asserted as a one-cycle pulse.
\end{enumerate}

\subsection{Free Flow}

\begin{enumerate}
  \item \texttt{rob\_last} pulse triggers Free Logic in the same cycle.
  \item Free Logic asserts \texttt{clear\_en} to the RTT Register File:
        the entry is invalidated.
  \item Free Logic simultaneously asserts \texttt{RD\_FREE\_VALID} with
        \texttt{RD\_FREE\_TAG} to the Tag Allocator, decrementing the
        occupancy counter for the freed AXID.
\end{enumerate}

\subsection{rob\_last Routing}

\texttt{rob\_last} is a sideband signal to Merge Logic. It travels in
parallel with, not through, the ROB. The ROB emits
\texttt{data[511:0]}, \texttt{status[3:0]}, and \texttt{AXID[5:0]} to
Merge Logic; the RTT emits \texttt{rob\_last} to Merge Logic
coincident with the final data beat. Merge Logic combines both to
assert \texttt{RLAST} on the AXI R channel. There is no feedback path
from the RTT into the ROB's internal logic.

% -----------------------------------------------------------------------
\section{Timing Notes}

\begin{itemize}
  \item Alloc Write Logic is registered: the register file write takes
        effect on the cycle following \texttt{TAG\_VALID}.
  \item Completion Update Logic is registered: \texttt{cmp\_pkts}
        increment takes effect on the cycle following \texttt{RESP\_VALID}.
  \item Compare Logic is combinational: \texttt{rob\_last} is available
        in the same cycle as the \texttt{cmp\_pkts} update is visible.
        In practice this means \texttt{rob\_last} appears one cycle
        after the final \texttt{RESP\_VALID}.
  \item Free Logic fires in the same cycle as \texttt{rob\_last}:
        the register file entry is cleared and \texttt{RD\_FREE\_VALID}
        is asserted without additional latency.
  \item Simultaneous allocation and completion on the same tag are not
        possible by construction — a tag is only allocated after its
        previous occupancy slot is freed, and a completion can only
        arrive for an in-flight tag.
\end{itemize}

% -----------------------------------------------------------------------
\section{Edge Cases}

\begin{itemize}
  \item \textbf{Single-beat transaction (ARLEN = 0):} \texttt{exp\_pkts = 1}.
        \texttt{rob\_last} is asserted on the first and only completion.
        Handled identically to multi-beat transactions.
  \item \textbf{Error completion:} \texttt{resp\_type == 2'b10} (ERR)
        for a read transaction is treated as a valid completion event.
        \texttt{cmp\_pkts} is incremented normally so the transaction
        progresses to \texttt{rob\_last}. Error status is stored in the
        ROB slot and forwarded to the R Response Generator as DECERR.
        The RTT does not need to distinguish error from normal
        completions.
  \item \textbf{WR\_ACK completions:} \texttt{resp\_type == 2'b01}
        completions are ignored by the RTT entirely. They are processed
        by the Response Tracker on the write path.
  \item \textbf{Over-completion protection:} \texttt{cmp\_pkts} must
        never exceed \texttt{exp\_pkts}. This is guaranteed by
        construction since MC Core returns exactly one completion per
        issued read packet. A design-time assertion should be included
        in the verification environment.
\end{itemize}

\end{document}
